% Draft: guidance, task settings, baselines. Numbers are from the runs listed in the comments; nothing was rounded by hand.
% TODO markers are where a number or a decision is still open.

\section{Planning by Guided Sampling}
\label{sec:guidance}

\subsection{Guidance on the initial noise}
\label{sec:guidance-method}

Our world model is an unconditional generative model over point-cloud trajectories: given the first observed frame,
it samples the future motion of every particle in the scene, including the particles of the robot's gripper.
We never condition the model on actions. Instead we plan by steering the sampler so that the generated future
satisfies a task objective, and then execute the gripper motion contained in that future.

Sampling is deterministic given the initial noise. We therefore treat the initial noise as the decision variable.
Let $\epsilon$ denote the noise tensor that the sampler consumes, and let $X(\epsilon) = (x_0, x_1, \dots, x_T)$ be the
generated trajectory, where $x_0$ is the observed first frame. For a task objective $J$ we solve
\begin{equation}
  \epsilon^\star = \arg\min_{\epsilon} \; J\big(X(\epsilon)\big)
\end{equation}
by differentiating through the full sampling chain and taking a small number of gradient steps on $\epsilon$.
Gradients are obtained by backpropagation with gradient checkpointing over the sampler's model calls.

\paragraph{Norm-preserving updates.}
A gradient step moves the noise off the region the model was trained on. After every step we rescale each
independent noise draw back to the norm it had when it was sampled (\emph{shell reprojection}):
\begin{equation}
  \epsilon_i \leftarrow \epsilon_i \, \frac{\lVert \epsilon_i^{(0)} \rVert}{\lVert \epsilon_i \rVert}.
\end{equation}
The guidance strength is then set entirely by the step size and the number of steps; there is no separate guidance weight.

\paragraph{Optimizers.}
We use one of two update rules.
\emph{Adam} on the noise, with a constant learning rate.
\emph{Unit-gradient} steps, $\epsilon_i \leftarrow \epsilon_i - \eta\sqrt{D}\, g_i/\lVert g_i\rVert$, where $D$ is the
dimension of draw $i$; each draw moves by a fixed fraction $\eta$ of its own expected norm per step, independent of
the gradient's scale.

\paragraph{Population.}
We optimize a population of $B$ noise samples of the same scene and the same target in parallel (the batch is the
population). The executed plan is the member with the lowest \emph{imagined} objective, i.e.\ the objective evaluated
on the model's own generated trajectory. No simulator or real-world information is used to choose the plan.

\subsection{Objectives}
\label{sec:objectives}

All objectives are differentiable functions of the generated trajectory. Let $c_t$ be the centroid of the object's
particles at frame $t$, $g$ the target, and $p_{t,j}$ the particles of the scene.

\begin{itemize}
  \item \textbf{Goal.} $J_{\text{goal}} = \lVert c_T - g \rVert$, the distance of the object's centroid at the last
        frame to the target. This is the only term in the benchmark tasks.
  \item \textbf{Obstacle clearance.} Each obstacle is an axis-aligned box. For every particle and every generated
        frame we compute the signed distance $d$ to the box and penalize the shortfall
        $s = \max(0,\, m - d)$ below a clearance margin $m$:
        $J_{\text{box}} = \operatorname{mean}_{t,j}(s) + \max_{t,j}(s)$.
        The margin acts as a buffer around the obstacle's volume.
  \item \textbf{Floor.} Mean depth of particles below the table plane (used only where noted).
  \item \textbf{Upright.} For placing tasks we penalize the tilt of the object at the end of the trajectory:
        the rotation $R_t$ of the object's particles relative to frame~0 is recovered by a differentiable Kabsch fit,
        and $J_{\text{up}} = \frac{1}{k}\sum_{t=T-k+1}^{T} \arccos\big((R_t)_{zz}\big)$ over the last $k = 5$ frames.
\end{itemize}
The total objective is $J = J_{\text{goal}} + w_{\text{box}}J_{\text{box}} + w_{\text{up}}J_{\text{up}}$.

\subsection{Tasks and settings}
\label{sec:task-settings}

Table~\ref{tab:guidance-settings} lists the settings per task. All timings are for warm, compiled (bf16) generation on
one NVIDIA L40S and exclude simulation.

\begin{table}[t]
\centering
\small
\caption{Guidance settings per task.}
\label{tab:guidance-settings}
\begin{tabular}{lccc}
\toprule
 & Pick-to-location & Moulding (dough) & Pick, carry and place (gallery) \\
\midrule
Model & pickup, 104.8M & mould, 104.3M & pickup, 104.8M \\
Horizon (transitions) & 60 (3.0\,s at 20\,fps) & 150 & 60 \\
Sampler & streaming, window 16 & streaming, keep 3, full churn & streaming, window 16 \\
Denoising steps / frame & 10 (Euler) & 6 & 10 (Euler) \\
Optimized noise & initial pose noise & bank draws at level $\ge 0.8$ & initial pose noise$^\dagger$ \\
Optimizer & Adam, lr 0.1 & unit-gradient, $\eta = 0.1$ & Adam, lr 0.1 \\
Iterations & 10 & 6 & 40 (80 with obstacles) \\
Shell reprojection & on & on & on \\
Population & 4 & 3 & $2 \times 4$ (two seeds) \\
Objective & goal & goal & goal (+ box, upright)$^\dagger$ \\
Plan time & $126.5 \pm 0.8$\,s & $113.8$\,s (mean) & $\approx 8$\,min per seed \\
\bottomrule
\end{tabular}
\\[2pt]\footnotesize $^\dagger$ Gallery only: see the comment in the pick-carry-place paragraph.
\end{table}

\paragraph{Pick-to-location.}
Five objects (milk carton, peanut butter jar, soda can, pasta box, water bottle) from held-out episodes, each with five
targets: target $=$ frame-0 object centroid $+ (\Delta x, \Delta y, 0.15\,\text{m})$, with $(\Delta x, \Delta y)$
uniform in a 0.15\,m disc, drawn once with a fixed seed and shared by all objects and all methods (25 problems).
Only frame~0 of the recorded episode is used.

\paragraph{Moulding.}
The dough-moulding task moves the centroid of a deformable dough by pushing it with a paddle. Fifty scenarios pin a
dataset episode and an in-plane target offset of 7.5--15\,cm drawn with a fixed seed. The sampler re-forms every
partially denoised row with fresh noise at every step (full churn, $\eta = 1$) on both the paddle and the dough
particles. Only the noise draws at high noise levels ($\ge 0.8$) are optimized.

\paragraph{Pick, carry and place.}
For the qualitative gallery we use harder variants of the lift task: air targets 30\,cm away and 15\,cm up, ground
targets behind a barrier (a cereal box standing on its long edge, 21\,cm tall, or a row of three cracker boxes, 15\,cm
tall), and placement on top of a 19.6\,cm cabinet with the upright term ($w_{\text{up}} = 0.3$). Obstacle cases use
a clearance margin of 25\,mm with $w_{\text{box}} = 10$; placement on the cabinet uses a 3\,mm margin. Every obstacle
case is paired with a no-obstacle run to the same target.
% GALLERY ONLY -- NOT used for any benchmark number (the 25 pick runs, cyl_t*_it10, were run 2026-09-23 without it):
% the gallery runs keep the pose noise of the first 24 frames fixed (freeze_frames 24) and add a pre-grasp stillness
% term (still_weight 5: mean object-centroid displacement from frame 0 over frames < 24). Decide whether the paper
% mentions this for the figure; it is not part of the method used for Tables pick/mould.
% TODO (paper choice): the gallery picks the member to render using the simulator (see Sec. evaluation).
% Say so explicitly in the figure caption, or restrict the figure to the imagined-objective pick.

\paragraph{Block pushing.}
% Settings from the push sweep (block pushing, SAPIEN): settled recipe.
Population 3, 6 denoising steps per frame, 6 unit-gradient iterations with $\eta = 0.1$, shell reprojection, and only
draws at noise levels $\ge 0.5$ optimized; about 20\,s per plan.
% TODO: final push numbers and whether rl_push is reported.

\subsection{Evaluation in simulation}
\label{sec:eval-sim}

The model plans in its own imagination; we evaluate by executing the planned gripper motion in a physics simulator
that the model has never seen.
\begin{itemize}
  \item \textbf{Pickup (Isaac Sim, floating hand).} A floating Franka hand tracks the gripper pose recovered from the
        generated gripper particles (rigid fit per frame) at 20\,fps. The gripper is binary: commanded closed when the
        generated finger opening falls below 35\,mm. The object is the convex hull of its frame-0 particles, with the
        density and friction of the source asset.
  \item \textbf{Pickup (Isaac Sim, commanded arm).} The same hand trajectory executed by a full, fixed-base Franka arm
        via damped-least-squares inverse kinematics with PD joint control and velocity feed-forward, started from the
        episode's recorded joint configuration. This adds reach and joint limits; obstacles are simulated as fixed
        collision bodies. The rendered figures show this arm.
  \item \textbf{Moulding (Genesis MPM).} The generated paddle poses are tracked by a PD-controlled paddle
        (25 simulation steps per model frame) in an MPM dough initialized from the frame-0 dough particles.
\end{itemize}
The headline metric is the distance from the simulated object centroid at the last frame to the target
(pickup), and the distance closed toward the target, $\lVert c_0 - g\rVert - \lVert c^{\text{sim}}_T - g\rVert$ (moulding).

\subsection{Baselines}
\label{sec:baselines}

All baselines receive exactly the same problems: the same frame-0 state, the same targets, the same horizon, and the
same scoring. Each chooses its plan by its own predicted objective, never by the simulator.

\paragraph{No guidance: sample-and-select MPC with the same model.}
To isolate the contribution of gradient guidance, we plan with the same world model and sampler but without
gradients. At every replanning step we draw $K$ fresh noise samples, generate $K$ continuations of the committed
trajectory prefix, keep the single candidate whose imagined final object centroid is closest to the target
(argmin, no averaging), commit its next 5 frames, and repeat until the horizon is reached. Committed frames enter the
sampler as clean context; every model call sees a full-size window, which slides back over committed frames near the
end rather than shrinking. $K$ is chosen so that one plan costs about as much compute as guidance:
$K = 20$ for pickup (138\,s) and $K = 6$ for moulding (151.5\,s).

\paragraph{PointWorld.}
PointWorld is an action-conditioned point-cloud world model. We train it on the same data and plan with the same
sample-and-select MPC: per replanning step, $K = 88$ candidate gripper trajectories drawn around the current plan
with Gaussian perturbations whose scale matches the recorded per-frame motion statistics (7.5\,mm translation,
0.06\,rad rotation), argmin selection on the predicted object-centroid distance to the target, and a horizon that
shrinks to the fixed final frame. $K$ is set to match our planning time (148.9\,s per plan).
% TODO: state PointWorld training steps / checkpoint used (pick benchmark used a 74k/200k checkpoint).

\paragraph{GNS.}
% TODO: GNS protocol and numbers. Same benchmark files (sweeps/pick_benchmark_5x5*.{json,npz}) were handed over.

\subsection{Results}
\label{sec:guidance-results}

\begin{table}[t]
\centering
\small
\caption{Pick-to-location, 25 problems. Distance of the simulated (Isaac, floating hand) object centroid to the target
at the last frame.}
\label{tab:pick-results}
\begin{tabular}{lcccc}
\toprule
Planner & median (mm) & $\le 50$\,mm & $\le 60$\,mm & plan time (s) \\
\midrule
Guided (ours) & \textbf{38} & \textbf{18/25} & \textbf{21/25} & $126.5 \pm 0.8$ \\
No guidance, MPC ($K = 20$) & 121 & 2/25 & 4/25 & $138.1 \pm 0.6$ \\
PointWorld, MPC ($K = 88$) & 190 & 0/25 & 0/25 & $148.9 \pm 20.8$ \\
GNS & -- & -- & -- & -- \\
\bottomrule
\end{tabular}
\end{table}

\begin{table}[t]
\centering
\small
\caption{Moulding, 50 scenarios (targets 7.5--15\,cm away). Distance closed toward the target in Genesis.}
\label{tab:mould-results}
\begin{tabular}{lccccc}
\toprule
Planner & mean (mm) & median (mm) & $\ge 20$\,mm & $\ge 30$\,mm & plan time (s) \\
\midrule
Guided (ours) & \textbf{26.8} & \textbf{26.0} & \textbf{40/50} & \textbf{15/50} & 113.8 \\
No guidance, MPC ($K = 6$) & 14.8 & 14.2 & 12/50 & 1/50 & $151.5 \pm 0.9$ \\
\bottomrule
\end{tabular}
\end{table}
% Mould std of distance closed (guided): 8.0 mm; MPC: 7.6 mm. Guided closes more than MPC on 43/50 paired scenarios.
% PointWorld predicted its own plans to 29.9 mm median on the pick benchmark; in Isaac the object was not lifted (peak lift 4.8 mm over all 25).
